\section{Interpretacion del caso}

\subsection{Analisis horizontal}

\begin{center}
\begin{tabular}{lrrr}
\toprule
Partida & 2024 & 2025 & Variacion \% \\
\midrule
Ventas netas & S/ 180,000.00 & S/ 215,000.00 & 19.44\% \\
Costo de ventas & S/ 108,000.00 & S/ 133,000.00 & 23.15\% \\
Utilidad bruta & S/ 72,000.00 & S/ 82,000.00 & 13.89\% \\
Utilidad neta & S/ 21,855.00 & S/ 25,027.00 & 14.51\% \\
Total activo & S/ 200,000.00 & S/ 225,000.00 & 12.50\% \\
Patrimonio & S/ 130,000.00 & S/ 143,000.00 & 10.00\% \\
\bottomrule
\end{tabular}
\end{center}

La empresa incrementa sus ventas en 19.44\%, lo cual es favorable. Sin embargo, el costo de ventas aumenta 23.15\%, es decir, crece mas rapido que las ventas. Esto reduce la eficiencia operativa y presiona el margen bruto.

\subsection{Analisis vertical}

\begin{center}
\begin{tabular}{lrr}
\toprule
Partida & 2024 & 2025 \\
\midrule
Ventas netas & 100.00\% & 100.00\% \\
Costo de ventas & 60.00\% & 61.86\% \\
Utilidad bruta & 40.00\% & 38.14\% \\
Utilidad neta & 12.14\% & 11.64\% \\
Activo corriente & 47.50\% & 50.22\% \\
Patrimonio & 65.00\% & 63.56\% \\
\bottomrule
\end{tabular}
\end{center}

El costo de ventas pasa de representar 60.00\% a 61.86\% de las ventas. En consecuencia, el margen bruto baja de 40.00\% a 38.14\%. La utilidad neta tambien reduce su peso relativo.

En el balance, el activo corriente gana participacion y el patrimonio reduce su peso relativo frente al total activo. Esto debe revisarse junto con ratios de liquidez y solvencia.

\subsection{Interpretacion gerencial del caso Empresa ABC}

El caso Empresa ABC muestra que el crecimiento de ventas no debe interpretarse de forma aislada. Aunque las ventas aumentan 19.44\%, el costo de ventas aumenta 23.15\%, lo que evidencia una presion sobre la eficiencia operativa.

Desde una perspectiva gerencial, la empresa podria revisar proveedores, precios, productividad, procesos internos y politicas de costos. Tambien podria configurar alertas en un dashboard financiero para detectar automaticamente cuando los costos crecen mas rapido que las ventas.

La lectura vertical confirma el hallazgo: el costo de ventas pasa de 60.00\% a 61.86\% de las ventas, reduciendo la utilidad bruta relativa. Esto indica que la empresa debe proteger sus margenes y no evaluar el crecimiento solo por el aumento de ingresos.
